\subsection{Out-of-Corpus Generalisation (WikiText-2)}
\label{sec:out_of_corpus}

\paragraph{Task Description.}
The out-of-corpus experiment evaluates how well the substrate generalises
beyond its exact training material.  The ingestion corpus is 10 samples
from the \texttt{wikitext-2-raw-v1} training split (processed Wikipedia
articles, mean length $\sim$350 tokens).  Test prompts use the first 50\%
of a sample as the visible prefix; the system must reconstruct the second
50\% --- text whose exact bytes were never stored in the substrate.

Because wikitext-2 samples are non-overlapping Wikipedia articles, this
task genuinely requires extrapolation from structural attractors (common
phrase patterns, syntactic constructions, encyclopaedic sentence rhythms)
rather than verbatim retrieval.

\paragraph{Results.}
Across $N = 10$ test samples the mean weighted score was 0.266.

\paragraph{Assessment.}
Partial generalisation was observed.  Common Wikipedia phrasing patterns
(article structure, parenthetical citations, passive voice constructions)
are recoverable, while domain-specific terminology and specific named
entities are not, as expected for a 10-sample ingestion corpus.

Figure~\ref{fig:out_of_corpus_map} shows the trial outcome map.
